
We opened this paper with a reversal: semantic entropy wins on TriviaQA by +0.155 AUROC, then loses on TruthfulQA by $-0.066$ AUROC. That reversal is not a failure of the method --- it is a diagnostic signal revealing the condition under which semantic-level uncertainty estimation provides genuine value.

This paper provides the first unified comparison of SE, TE, SCG, and VC at Llama-2-7B scale, with direct mechanism measurement and cross-benchmark testing. We establish that SE substantially outperforms TE on TriviaQA (AUROC 0.717 vs. 0.562, gap = +0.155, gap CI excludes zero), confirming the SE > TE ordering at sub-65B scale. We identify the mechanism: intra-cluster TE variance = 7.152 $nats^2$ ($71\times$ threshold) across 76/98 questions confirms TE aggregates paraphrase noise that SE's NLI clustering removes. We identify the task-structure scope condition: the ordering reverses on TruthfulQA (TE = 0.511 > SE = 0.445), where deterministic-wrong outputs defeat SE's filtering advantage --- though NLI model size is a confound that future work must disentangle. We characterize two alternative method failure modes: BERTScore SCG diverges from SE by 0.336 AUROC on short QA (3--$11\times$ the 0.03 equivalence gate), and VC at 7B scale produces degenerate confidence (ECE = 0.430, 5 distinct values, ~60\% at 95\%).

\textbf{Future directions.} (1) Re-run H-C1 with \texttt{nli-deberta-v3-large} to disentangle NLI model size from task structure in the TruthfulQA reversal. (2) Run N=500 extension of H-E1 to confirm gap stability at larger sample size. (3) Test SelfCheckNLI (NLI-based SCG variant) to determine whether NLI-based consistency recovers SE-equivalence on short QA. (4) Scale to Llama-13B and 70B to test whether the SE > TE gap grows with scale, as suggested by Kuhn et al. [2023]'s 65B results.

The uncertainty method that wins on TriviaQA loses on TruthfulQA. Rather than choosing between benchmarks, practitioners should ask what each task's incorrect outputs look like --- diverse or deterministic --- and select accordingly. Understanding this scope condition is more valuable than an unconditional method ranking, and provides a foundation for principled uncertainty method selection as deployment settings grow more varied.

